% TikZ 图：等量同种电荷的电场线
% Day 15：电荷、电场与电场强度
\begin{tikzpicture}[scale=0.9,>=Latex]
  \coordinate (A) at (-1.4,0);
  \coordinate (B) at (1.4,0);

  \filldraw[red] (A) circle (4pt);
  \node[below,fill=white,inner sep=1pt] at (A) {$+Q$};
  \filldraw[red] (B) circle (4pt);
  \node[below,fill=white,inner sep=1pt] at (B) {$+Q$};

  % 同种正电荷的电场线都从电荷出发，向远处延伸。
  % 面向彼此的一侧会向上下分流；电场线不连接、不相交。
  \draw[->,thick] (-1.58,0.00) -- (-3.10,0.00);
  \draw[->,thick] (-1.53,0.15) .. controls (-1.95,0.35) and (-2.45,0.75) .. (-3.05,1.25);
  \draw[->,thick] (-1.53,-0.15) .. controls (-1.95,-0.35) and (-2.45,-0.75) .. (-3.05,-1.25);
  \draw[->,thick] (-1.30,0.18) .. controls (-1.05,0.55) and (-0.88,1.15) .. (-0.82,1.95);
  \draw[->,thick] (-1.30,-0.18) .. controls (-1.05,-0.55) and (-0.88,-1.15) .. (-0.82,-1.95);

  \draw[->,thick] (1.58,0.00) -- (3.10,0.00);
  \draw[->,thick] (1.53,0.15) .. controls (1.95,0.35) and (2.45,0.75) .. (3.05,1.25);
  \draw[->,thick] (1.53,-0.15) .. controls (1.95,-0.35) and (2.45,-0.75) .. (3.05,-1.25);
  \draw[->,thick] (1.30,0.18) .. controls (1.05,0.55) and (0.88,1.15) .. (0.82,1.95);
  \draw[->,thick] (1.30,-0.18) .. controls (1.05,-0.55) and (0.88,-1.15) .. (0.82,-1.95);

  \filldraw[black] (0,0) circle (1.6pt);
  \node[below,fill=white,inner sep=1pt] at (0,0) {$O$};
  \node[font=\small] at (0,-2.20) {中点处两场强等大反向，合场强为零};
\end{tikzpicture}
